The orbit of Phobos and the rotation of the planet are in the same direction,
so the angular velocity of the moon relative to Mars, $\omega_{P/M}$, is equal
to the difference
\[
  \omega_{P/M} = \omega_M - \omega_P
\]
\hfill[2]

where $\omega_M$ and $\omega_P$ are the angular rotation velocity of Mars and
the angular orbital velocity of Phobos respectively.
For $\omega_M$ we have $\omega_M = 2\pi/T_M$ where $T_M$ is the period \hfill[1]

The time above the horizon, $t$, is found from the relation (see drawing):
\[
  t = 2\alpha/\omega_{P/M}
\]
\hfill[2]

where $\alpha = R_M/R_P$.
% sic: source omits the cosine; the geometry (and the numerical value below) give cos(alpha) = R_M/R_P
\hfill[1]

% FIGURE NEEDED: two concentric circles (Mars of radius R_M and the orbit of Phobos of radius R_P) with the horizon plane drawn tangent to Mars at the sub-observer point, radii R_M and R_P marked to the intersection of the orbit with the horizon (2011 TH short solutions, p. 73)

Find the angular velocity of Phobos $\omega_P$ assuming gravity is equal to
the centripetal force:
\[
  \omega_P = \mathrm{sqrt}(\,GM_M/R^{3/2}_{\;P})
  % sic: source prints sqrt(GM_M / R_P^{3/2}); dimensionally it should be sqrt(GM_M / R_P^3)
\]
\hfill[2]

Substituting:
\begin{align*}
  \omega_M &= 7.0882\cdot 10^{-5}\,\mathrm{s}^{-1}, \\
  \omega_P &= 2.2784\cdot 10^{-4}\,\mathrm{s}^{-1}, \\
  \omega_{P/M} &= -1.5696\cdot 10^{-4}\,\mathrm{s}^{-1}
\end{align*}
where ``$-$'' means that the moon ``leads'' the rotation of the planet;
substitute the absolute value into the equation,

thus
$\alpha = 68^\circ\!.794 = 1.2007$\,rad.
and $t \approx 15300\,\mathrm{s} = 4h\,15m$. \hfill[2]
